\chapter{RESUMEN DEL PROYECTO}

\section{Contexto y justificación}

Junto al aumento exponencial de sensores y dispositivos conectados a internet, existen retos lejos de ser solventados causados por la gran variedad de protocolos, formatos y estándares de los fabricantes de estos dispositivos. El volumen de datos generado por el Internet de las Cosas (IoT) impone requisitos que deben ser satisfechos por arquitecturas capaces de procesar y analizar datos en tiempo real. Sobre la arquitectura Kappa se encuentran multitud de publicaciones, tutoriales y documentación de todo calibre que no llegan a ofrecer una implementación extremo a extremo de esta arquitectura para procesar y analizar datos por flujos. La mayoría de los recursos formativos y de referencia dependen en algún punto de servicios gestionados o plataformas propietarias, dificultando una experiencia de aprendizaje completa basada en herramientas de código abierto.

\section{Planteamiento del problema}

El problema que aborda el presente proyecto es la falta de una implementación de referencia, de código abierto y containerizada, de una arquitectura Kappa para IoT que integre bridging MQTT-Kafka, gobernanza de esquemas mediante Avro/Apicurio, y una estrategia de doble sumidero para consumo operacional y analítico. Por tanto, se propone contribuir a llenar vacíos de información detectados durante la investigación del estado del arte sobre arquitecturas públicas y de código abierto con un proyecto que tiene carácter académico y científico-técnico; no responde a un encargo empresarial y su aporte de innovación está en integrar en una sola arquitectura todos los elementos antes mencionados.

\section{Objetivos del proyecto}

El objetivo general es diseñar, desarrollar y validar un sistema escalable de microservicios para la ingesta, el procesamiento y el análisis operacional de datos IoT. Para cubrir ese alcance se plantean las tareas de analizar el estado del arte y las herramientas de código abierto disponibles, implementar un simulador de dispositivos IoT, implementar la arquitectura de microservicios containerizada; con servicios de ingesta sobre protocolo estándar, de gobernanza y validación de esquemas, de agregación y transformación, y de almacenamiento para series temporales y para consumo analítico, e implementar los servicios de visualización interactiva. Cada tarea se concreta en uno de cinco objetivos: persistencia de telemetría rápida y fiable, gobernanza del esquema de datos, procesamiento en flujo con baja latencia, visualización diferenciada operacional y analítica, y resiliencia del sistema, cada uno acompañado de un indicador clave de rendimiento (KPI) verificable.

\section{Resultados obtenidos}

Los cinco objetivos se cumplieron dentro de los umbrales fijados por sus KPIs. La telemetría validada se persiste dentro del tiempo objetivo y con una tasa de pérdida nula; la totalidad de los eventos ingeridos se valida contra el esquema Avro registrado sin fallos de deserialización; cada micro-lote de procesamiento se resuelve por debajo del límite de latencia sin incremento registrado; la visualización se ofrece mediante tres \textit{dashboards} de Grafana y tres páginas de Power BI que cubren el consumo energético, la eficiencia operativa, la detección de anomalías y el análisis histórico; y el sistema se recupera de la caída de sus servicios por debajo del tiempo objetivo y sin pérdida de datos. El resultado es la implementación completa, containerizada y reproducible de una arquitectura Kappa para IoT construida con herramientas de código abierto, disponible para su adopción y adaptación en proyectos similares de monitorización energética.

\section{Estructura de la memoria}

\textbf{Antecedentes y estado del arte} sitúa el crecimiento de los datos IoT, compara las arquitecturas Lambda y Kappa, revisa el ecosistema de herramientas de código abierto por capas y algunos retos abiertos, y formula el planteamiento del problema. \textbf{Objetivos} define el alcance, las tareas a desarrollar y los objetivos con sus KPIs y acciones. \textbf{Ética y compromiso social} recoge la relación del trabajo con la igualdad de género y con los Objetivos de Desarrollo Sostenible. \textbf{Desarrollo del proyecto} describe la planificación, la arquitectura y las herramientas empleadas, la preparación de los datos, cada componente del \textit{pipeline}, los recursos y el presupuesto, y presenta los resultados obtenidos para cada objetivo. \textbf{Ejecución y validación del sistema} documenta los modos de arrancar el \textit{pipeline} y las herramientas para medir KPIs y ejecutar pruebas de fallos. \textbf{Discusión} expone las limitaciones del proyecto y el impacto de los resultados. \textbf{Conclusiones} valora el cumplimiento de los objetivos y sus KPIs respectivos. \textbf{Futuras líneas de trabajo} plantea la evolución de la orquestación del sistema. Se culmina la memoria con la \textbf{bibliografía} y los \textbf{anexos}.
